% ============================================================
% Section 8: Limitations
% ============================================================

\section{Limitations}
\label{sec:limitations}

We identify eight limitations that bound the scope and applicability of our findings.

\paragraph{Open-weight models only.}
All experiments require direct access to \kvcache tensors extracted during inference.
This restricts applicability to open-weight models (Qwen, Llama, Mistral, and similar).
Closed-model APIs (GPT-4, Claude, Gemini) do not expose key-value states, and our methodology cannot be applied to them without provider cooperation.
Whether providers could implement cache-based monitoring internally is an open question, but external auditors cannot.

\paragraph{System prompt length sensitivity.}
When system prompts differ in token count between experimental conditions, the length difference leaks into encoding-phase features and inflates detection \auroc.
Phase~2b demonstrated that a 14--16 token difference (tokenizer-dependent) between honest and confabulation system prompts was sufficient to produce spuriously high discrimination.
Double residualization---controlling separately for system prompt tokens and generated tokens---is necessary under these conditions but reduces some comparisons to chance.
Equal-length system prompt padding eliminates this confound but constrains experimental design.
Any deployment using cache geometry for detection must ensure that prompt-length variation is controlled or residualized.

\paragraph{Architecture-specific calibration.}
No universal feature set works across all tested architectures.
Each architecture relies on different dominant features for cognitive mode discrimination: Qwen on key entropy, Llama on head norm standard deviation, Mistral on top singular value ratio.
Only three of eight features are universal in sign direction across architectures.
Practical deployment requires per-model geometric profiling---a calibration step that adds setup cost and limits scalability to novel architectures without prior characterization.

\paragraph{Small sample sizes on specific findings.}
Several findings rest on limited trial counts.
Jailbreak detection was tested on only two models.
The 70B same-prompt deception experiment (Experiment~44) comprised 20 trials per condition.
While effect sizes are large in these cases (\auroc $> 0.87$ for jailbreak), the confidence intervals are wide, and these findings require larger-scale replication before they should inform deployment decisions.
The original sycophancy experiment (Experiment~39, 16 trials) was falsified by input-length audit and replaced with a 600-trial replication (Section~\ref{sec:sycophancy}), resolving one of the most significant sample-size concerns.

\paragraph{Greedy decoding only.}
All detection experiments use greedy decoding (temperature $= 0$).
The effect of sampling temperature on cache geometry is explored only in the concordance study's Phase~B ICC analysis, which examines test-retest reliability under identical prompts, not the systematic effect of temperature variation on geometric features.
Deployed models typically use $\text{temperature} > 0$; whether stochastic sampling introduces noise that degrades detection accuracy, or whether the geometric signal is robust to sampling variation, is untested in the detection context.

\paragraph{Forced-instruction design.}
The majority of our experiments induce cognitive modes through explicit system prompt instructions (e.g., ``You must answer deceptively'' or ``Fabricate a plausible response'').
Whether the same geometric signatures appear in naturally occurring, uninstructed misalignment---a model that develops deceptive behavior through training dynamics rather than instruction---is only partially tested.
Phase~3c examined naturally generated confabulation (asking models questions they cannot answer correctly) and found marginal detection after \fwl control (\auroc 0.67--0.70), substantially below the instructed condition.
This gap between instructed and natural conditions is a significant concern for real-world applicability, where misaligned models will not have been told to be deceptive.

\paragraph{No per-comparison confidence intervals.}
Bootstrap confidence intervals are not reported for individual \auroc values or effect sizes; this limits the precision of our claims, particularly for comparisons with small sample sizes (e.g., sycophancy at $N = 16$ per condition).
The red-team battery reports aggregate bootstrap CIs that exclude 0.50, but per-comparison CIs would provide more informative precision bounds.
A partial replication ($N = 100$ per condition, 3 models, full input+output \fwl, BCa bootstrap from 10{,}000 resamples) has been completed for refusal-related comparisons.
The replication confirms that refusal detection \auroc values decrease from initial estimates (e.g., Qwen harmful-vs-benign: $0.788 \rightarrow 0.624$) but remain above chance for Qwen and Mistral.
Llama refusal detection collapses to chance ($0.551$, CI includes 0.50), confirming it as a length artifact.
Mistral shows the strongest confound-free signal ($0.787$, output AUROC = 0.505) because it complies with harmful prompts rather than refusing them.
Full per-comparison BCa CIs for all experimental comparisons remain a planned extension.

\paragraph{Causal gap.}
All evidence is correlational.
We have established that \kvcache geometry co-occurs with cognitive mode, discriminates between modes, and reveals structural relationships among modes.
We have not demonstrated that the geometry is causally necessary for the behavior or that intervening on the geometry changes the behavior.
Until the abliteration experiments outlined in Section~\ref{sec:discussion} are completed---removing cognitive-mode-associated directions from $W_K$ and testing for behavioral change---the possibility remains that cache geometry is an epiphenomenal byproduct of the actual computational mechanism rather than a participant in it.